% ============ 14 曲线的上同调 (原书 pp.89-104) ============
\chapter{曲线的上同调}\label{sec:14}
\index{曲线}
\index{Picard 群}
\index{紧支集上同调}
\index{Poincar\'e 对偶}
\index{Weil 配对}
\index{Hochschild-Serre 谱序列}
\index{可构造层}
\index{局部 Ext 群}
\index{有限域上的曲线}

本节研究曲线的平展上同调.这不仅是一般簇的平展上同调研究的范例，而且一般情形的许多证明都是从曲线的情形出发对簇的维数作归纳.

\subsection*{曲线的 Picard 群}

我们在上一节看到，对代数闭域 $k$ 上的连通非奇异曲线 $U$，
\[
H^{r}(U_{\mathrm{et}}, \mathbb{G}_{m}) = \begin{cases}
k^{\times}, & r = 0;\\
\operatorname{Pic}(U), & r = 1;\\
0, & r > 1.
\end{cases}
\]
我们希望用 Kummer 序列计算 $\mu_n$ 的上同调，但首先需要知道 $\operatorname{Pic}(U)$ 的结构.

$U$ 的 Picard 群可以由正合序列
\[
K^{\times} \to \bigoplus_{x \in U} \mathbb{Z} \to \operatorname{Pic}(U) \to 0
\]
定义.这里 $K = k(U)$，即 $U$ 上的有理函数域，求和遍历 $U$ 的所有闭点.对 $U$ 的闭点 $x$，设 $[x]$ 是对应于它的因子.于是 $\mathcal{D}iv(U)$ 的任何元素都可以写成有限和 $D = \sum_{x \in U} n_x [x]$，$n_x \in \mathbb{Z}$.$D = \sum n_x [x]$ 的次数是 $\sum n_x$.

现在设 $X$ 是代数闭域上的完备连通非奇异曲线.$X$ 上的非零有理函数 $f$ 的因子 $\mathrm{div}(f)$ 的次数为零(有理函数的极点与零点按重数计一样多).设 $\mathcal{D}iv^{0}(X)$ 表示次数为 0 的因子群，$\operatorname{Pic}^{0}(X)$ 是 $\mathcal{D}iv^{0}(X)$ 除以主因子子群的商.

\begin{proposition}\label{prop:14.1}
序列
\[
0 \to \operatorname{Pic}^{0}(X) \to \operatorname{Pic}(X) \to \mathbb{Z} \to 0
\]
正合.对与 $k$ 的特征互素的任何整数 $n$，
\[
z \mapsto nz \colon \operatorname{Pic}^{0}(X) \to \operatorname{Pic}^{0}(X)
\]
满射，其核等于秩 $2g$ 的自由 $\mathbb{Z}/n\mathbb{Z}$-模，其中 $g$ 是 $X$ 的亏格.
\end{proposition}

\begin{proof}
该序列是映射对 $K^{\times} \xrightarrow{\mathrm{div}} \mathcal{D}iv^{0}(X) \hookrightarrow \mathcal{D}iv(X)$ 的核-余核序列\footnote{映射对 $A \xrightarrow{f} B \xrightarrow{g} C$ 的核-余核正合序列是
\[
0 \to \operatorname{Ker}(f) \to \operatorname{Ker}(g \circ f) \to \operatorname{Ker}(g) \to \operatorname{Coker}(f) \to \operatorname{Coker}(g \circ f) \to \operatorname{Coker}(g) \to 0
\]
(CFT II 4.2).}的一部分.

第二个论断的证明更困难.首先假定 $k = \mathbb{C}$.选取 Riemann 曲面 $X(\mathbb{C})$ 上全纯微分的一组基 $\omega_1, \ldots, \omega_g$ 与 $H_1(X(\mathbb{C}), \mathbb{Z})$ 的一组基 $\gamma_1, \ldots, \gamma_{2g}$.设 $\Lambda$ 是由向量
\[
\left( \int_{\gamma_i} \omega_1, \ldots, \int_{\gamma_i} \omega_g \right), \qquad i = 1, \ldots, 2g
\]
生成的 $\mathbb{C}^{g}$ 的子群.对每对点 $z_0, z_1 \in X(\mathbb{C})$，选取从 $z_0$ 到 $z_1$ 的道路 $\gamma(z_0, z_1)$，设
\[
I(z_0, z_1) = \left( \int_{\gamma(z_0, z_1)} \omega_1, \ldots, \int_{\gamma(z_0, z_1)} \omega_g \right) \in \mathbb{C}^{g}.
\]
它在 $\mathbb{C}^{g}/\Lambda$ 中的像不依赖于道路 $\gamma(z_0, z_1)$ 的选取，而映射 $[z_1] - [z_0] \mapsto I(z_0, z_1)$ 通过线性延拓给出同态
\[
i \colon \mathcal{D}iv^{0}(X) \to \mathbb{C}^{g}/\Lambda.
\]
著名的 Abel 定理(Fulton, W., \emph{Algebraic Topology}, Springer 1995, 21.18)说 $i(D) = 0$，当且仅当 $D$ 主；同样著名的 Jacobi 反演定理说 $i$ 满(同上，21.32).因此，$i$ 诱导同构
\[
\operatorname{Pic}^{0}(X) \xrightarrow{\sim} \mathbb{C}^{g}/\Lambda.
\]
显然，对任何整数 $n$，$x \mapsto nx \colon \mathbb{C}^{g}/\Lambda \to \mathbb{C}^{g}/\Lambda$ 满射，且核为
\[
\frac{1}{n}\Lambda/\Lambda \cong \left(\frac{1}{n}\mathbb{Z}/\mathbb{Z}\right)^{2g} \cong (\mathbb{Z}/n\mathbb{Z})^{2g}.
\]
这完成了 $k = \mathbb{C}$ 情形的证明.

对任意代数闭域 $k$，证明类似，但需要 Jacobi 簇的代数理论.简言之，存在维数 $g$ 的阿贝尔簇 $J$，称为 $X$ 的 Jacobi 簇，使得 $\operatorname{Pic}^{0}(X) \cong J(k)$.对代数闭域 $k$ 上维数 $g$ 的任意阿贝尔簇 $J$ 与 $k$ 的特征互素的整数 $n$，$n \colon J(k) \to J(k)$ 满射，且核是秩 $2g$ 的自由 $\mathbb{Z}/n\mathbb{Z}$-模(见例如我的文章，\emph{Arithmetic Geometry}, Springer, 1986).
\end{proof}

\subsection*{\texorpdfstring{$\mu_n$ 的上同调}{mu_n 的上同调}}

\begin{proposition}\label{prop:14.2}
设 $X$ 是代数闭域 $k$ 上的完备连通非奇异曲线.对与 $k$ 的特征互素的任何 $n$，
\[
H^{0}(X, \mu_n) = \mu_n(k), \qquad H^{1}(X, \mu_n) \cong (\mathbb{Z}/n\mathbb{Z})^{2g}, \qquad H^{2}(X, \mu_n) \cong \mathbb{Z}/n\mathbb{Z},
\]
且 $r > 2$ 时 $H^{r}(X, \mu_n) = 0$.
\end{proposition}

\begin{proof}
由 Kummer 序列 (\ref{ex:7.9})，得到正合序列
\[
\cdots \to H^{r}(X_{\mathrm{et}}, \mu_n) \to H^{r}(X_{\mathrm{et}}, \mathbb{G}_{m}) \xrightarrow{n} H^{r}(X_{\mathrm{et}}, \mathbb{G}_{m}) \to \cdots.
\]
于是，论断可以由 (\ref{prop:14.1}) 与 (\ref{thm:13.7}) 读出.
\end{proof}

\begin{proposition}\label{prop:14.3}
设 $U$ 是代数闭域 $k$ 上的非奇异曲线.对与 $k$ 的特征互素的任何 $n$ 与 $U$ 的闭点 $x$，
\[
H^{2}_{x}(U, \mu_n) \cong \mathbb{Z}/n\mathbb{Z}; \qquad H^{r}_{x}(U, \mu_n) = 0 \quad (r \neq 2).
\]
\end{proposition}

\begin{proof}
设 $R$ 是 $\mathcal{O}_{U, u}$ 的 Hensel 化.由切除 (\ref{cor:9.8})，
\[
H^{r}_{x}(U, \mu_n) \cong H^{r}_{x}(V, \mu_n), \qquad V = \operatorname{Spec} R.
\]
现在对 $r > 0$，$H^{r}(V, \mathbb{G}_{m}) = 0$(定理~\ref{thm:13.7} 对 $V$ 也适用)，因此空间对 $(V, V \setminus x)$ 的正合序列
\[
\cdots \to H^{r}_{x}(V, \mathbb{G}_{m}) \to H^{r}(V, \mathbb{G}_{m}) \to H^{r}(V \setminus x, \mathbb{G}_{m}) \to \cdots
\]
为我们提供对所有 $r \neq 1$ 的同构
\[
H^{r-1}(V \setminus x, \mathbb{G}_{m}) \to H^{r}_{x}(V, \mathbb{G}_{m}).
\]
但 $V \setminus x = \operatorname{Spec} K$，其中 $K$ 是 $R$ 的分式域，而 Lang 定理~(\ref{item:13.5}c) 与命题~\ref{prop:13.6} 表明对 $r \geq 1$，$H^{r}(K, \mathbb{G}_{m}) = 0$.因此
\[
H^{1}_{x}(V, \mathbb{G}_{m}) \cong H^{0}(K, \mathbb{G}_{m})/H^{0}(V, \mathbb{G}_{m}) \cong \mathbb{Z}, \qquad H^{r}_{x}(V, \mathbb{G}_{m}) = 0 \quad (r \neq 1).
\]
命题现在由 Kummer 序列的正合序列
\[
\cdots \to H^{r}_{x}(V, \mu_n) \to H^{r}_{x}(V, \mathbb{G}_{m}) \xrightarrow{n} H^{r}_{x}(V, \mathbb{G}_{m}) \to \cdots
\]
推出.
\end{proof}

\begin{remark}\label{rem:14.4}
设 $M$ 是秩 1 的自由 $\mathbb{Z}/n\mathbb{Z}$-模，$M$ 也表示由 $M$ 定义的簇(或概形) $Y$ 上的常值层.则
\[
H^{r}(Y_{\mathrm{et}}, M) \cong H^{r}(Y_{\mathrm{et}}, \mathbb{Z}/n\mathbb{Z}) \otimes M \cong H^{r}(Y_{\mathrm{et}}, \mathbb{Z}/n\mathbb{Z}).
\]
当 $Y$ 是含 $n$ 个互异单位根的域 $k$ 上的簇时，该评注适用于 $M = \mu_n$.
\end{remark}

\subsection*{曲线的紧支集上同调}

首先回忆一些定义.域 $k$ 上的单变量函数域是 $k$ 上超越次数为 1 且在 $k$ 上有限生成的域 $K \supset k$，即 $K$ 是 $k(T)$($T$ 在 $k$ 上超越)的有限扩张.$k$ 上的曲线 $U$ 是正则的，如果对所有 $u$，$\mathcal{O}_{U, u}$ 是离散赋值环.非奇异曲线正则，但正则曲线不必非奇异，除非 $k$ 完全(按定义，域 $k$ 上的簇 $U$ 非奇异，如果 $U_{k^{\mathrm{al}}}$ 在 $k^{\mathrm{al}}$ 上非奇异---用概形语言，$U$ 在 $\operatorname{Spec} k$ 上光滑).

设 $U$ 是域 $k$ 上的连通正则曲线.则 $K \overset{\mathrm{def}}{=} k(U)$ 是 $k$ 上的单变量函数域，而映射 $u \mapsto \mathcal{O}_{U, u}$ 是从 $U$ 到满足如下条件的离散赋值环 $R \subset K$ 的集合的单射：
\begin{quote}
($\ast$) $R$ 含 $k$ 且以 $K$ 为分式域.
\end{quote}
单射性的证明使用簇分离(见``记号与约定'')---对``原点翻倍的仿射线''(AG 4.10)，该映射不是单射.

反过来，每个单变量函数域 $K$ 都典范地联系一条连通完备正则曲线 $X$.事实上，可以把 $X$ 定义为满足 ($\ast$) 的 $K$ 中所有离散赋值环的集合，并赋予其真闭子集为有限集的拓扑.对 $X$ 中的每个开集 $U$，定义
\[
\Gamma(U, \mathcal{O}_X) = \bigcap_{R \in U} R.
\]
则 $(X, \mathcal{O}_X)$ 是以 $K$ 为有理函数域的完备曲线，$X$ 在点 $R$ 处的局部环就是 $R$ 本身.(见 Hartshorne 1977, I.6.)

设 $U$ 是 $k$ 上的连通正则曲线，$X$ 是与 $k(U)$ 典范联系的连通完备正则曲线.则映射
\[
j \colon U \to X, \qquad j(u) = \mathcal{O}_{U, u},
\]
把 $U$ 实现为 $X$ 的开子簇.对 $U$ 上的层 $F$，定义
\[
H^{r}_{c}(U, F) = H^{r}(X, j_{!}F),
\]
并把 $H^{r}_{c}(U, F)$ 称为 $F$ 的紧支集上同调群.关于这个术语的解释见下面 \S 18.

因为 $j_{!}$ 正合，$U$ 上层的短正合序列给出上同调群的长正合序列.然而，因为 $j_{!}$ 不保持内射对象，$H^{r}_{c}(U, F)$ 不是 $H^{0}_{c}(U, \cdot)$ 的第 $r$ 个右导出函子.

\begin{proposition}\label{prop:14.5}
对代数闭域 $k$ 上的任何连通正则曲线 $U$ 与 $k$ 的特征不能整除的整数 $n$，有典范同构
\[
H^{2}_{c}(U, \mu_n) \to \mathbb{Z}/n\mathbb{Z}.
\]
\end{proposition}

\begin{proof}
设 $j \colon U \hookrightarrow X$ 是 $U$ 到完备正则曲线的典范包含，$i \colon Z \hookrightarrow X$ 是 $U$ 在 $X$ 中的补集.把 $\mu_n$ 看作 $X$ 上的层.由序列(见命题~\ref{prop:8.15})
\[
0 \to j_{!}j^{*}\mu_n \to \mu_n \to i_{*}i^{*}\mu_n \to 0
\]
得到正合序列
\[
\cdots \to H^{r}_{c}(U, \mu_n) \to H^{r}(X, \mu_n) \to H^{r}(X, i_{*}i^{*}\mu_n) \to \cdots.
\]
但
\[
H^{r}(X, i_{*}i^{*}\mu_n) \cong H^{r}(Z, i^{*}\mu_n) = 0 \quad (r > 0),
\]
从而
\[
H^{2}_{c}(U, \mu_n) \cong H^{2}(X, \mu_n) \cong \mathbb{Z}/n\mathbb{Z}.
\]
\end{proof}

\begin{remark}\label{rem:14.6}
对任何 $x \in U$，$U$ 上的层 $F$ 与 $r \geq 0$，有典范映射 $H^{r}_{x}(U, F) \to H^{r}_{c}(U, F)$.对 $F = \mu_n$ 与 $r = 2$，该映射与 (\ref{prop:14.3}) 和 (\ref{prop:14.5}) 中的同构相容.
\end{remark}

\subsection*{Poincar\'e 对偶定理}

本小节中，$U$ 是代数闭域 $k$ 上的连通正则曲线，$n$ 是 $k$ 的特征不能整除的整数.

所谓有限局部常值层 $F$，我指被 $n$ 消灭，茎有限的局部常值层.于是，对某个有限平展覆盖 $U' \to U$，$F|_{U'}$ 是由有限 $\mathbb{Z}/n\mathbb{Z}$-模 $M$ 定义的常值层，而给出 $U$ 上的有限局部常值层 $F$ 就是给出带 $\pi_1(U, \bar{u})$ 连续作用的有限 $\mathbb{Z}/n\mathbb{Z}$-模.(见 \S 6.)

\begin{theorem}\label{thm:14.7}[Poincar\'e 对偶]
对 $U$ 上的任何有限局部常值层 $F$ 与整数 $r \geq 0$，有有限群的典范完美配对
\[
H^{r}_{c}(U, F) \times H^{2-r}(U, \mathcal{L}F(1)) \to H^{2}_{c}(U, \mu_n) \cong \mathbb{Z}/n\mathbb{Z}.
\]
\end{theorem}

配对 $M \times N \to C$ 称为完美的，如果诱导映射
\[
M \to \operatorname{Hom}(N, C), \qquad N \to \operatorname{Hom}(M, C)
\]
都是同构.层 $\mathcal{L}F(1)$ 是
\[
V \mapsto \operatorname{Hom}_{V}(F|_{V}, \mu_n|_{V}).
\]
注意，若 $G = \mathcal{L}F(1)$，则 $F = \mathcal{L}G(1)$.

关于上述定理与拓扑空间通常的 Poincar\'e 对偶定理的关系的讨论见 \S 24.

我推迟到后面描述这个配对.在定理证明的梗概中，我省略图交换性的证明.然而，必须理解定理讲的是特定的配对，即它不只是说这些群互为对偶，而是说它们带特定的配对互为对偶.

\begin{example}\label{ex:14.8}
设 $X$ 是代数闭域 $k$ 上的完备连通非奇异曲线.在这种情况下，配对是
\[
H^{r}(X, F) \times H^{2-r}(X, \mathcal{L}F(1)) \to H^{2}(X, \mu_n) \cong \mathbb{Z}/n\mathbb{Z}.
\]
取 $F = \mathbb{Z}/n\mathbb{Z}$ 并把两个群之一与 $\mu_n$ 作张量积(比较 14.4)，得到配对
\[
H^{r}(X, \mu_n) \times H^{2-r}(X, \mu_n) \to H^{2}(X, \mu_n \otimes \mu_n) \cong \mu_n.
\]
对 $r = 1$，它可以等同于 $X$ 的 Jacobi 簇上 $n$ 阶点上的 Weil 配对
\[
\operatorname{Jac}(X)_n \times \operatorname{Jac}(X)_n \to \mu_n
\]
(AV I \S 13).
\end{example}

\subsection*{Poincar\'e 对偶证明的梗概}

在证明中，我们经常以如下形式使用五引理.考虑行正合的交换图
\[
\begin{tikzcd}
0 \arrow[r] & A' \arrow[r] \arrow[d, "\alpha"] & B' \arrow[r] \arrow[d, "\beta"] & C' \arrow[r] \arrow[d, "\gamma"] & D' \arrow[d, "\delta"] \\
0 \arrow[r] & A \arrow[r] & B \arrow[r] & C \arrow[r] & D
\end{tikzcd}
\]
则
\[
\beta, \delta \text{ 单射}, \alpha \text{ 满射} \Rightarrow \gamma \text{ 单射};
\qquad
\alpha, \gamma \text{ 满射}, \delta \text{ 单射} \Rightarrow \beta \text{ 满射}.
\]

第 0 步：对任何挠层 $F$，群 $H^{r}(U, F)$ 与 $H^{r}_{c}(U, F)$ 对 $r > 2$ 为零.

\begin{proof}
关于 $H^{r}(U, F) = 0$ 的消失见下面的定理~\ref{thm:15.1}.它蕴含对 $r > 2$，$H^{r}_{c}(U, F) \overset{\mathrm{def}}{=} H^{r}(X, j_{!}F) = 0$.
\end{proof}

于是，定理对 $r \neq 0, 1, 2$ 成立.

为方便，写 $T^{r}(U, F) = H^{2-r}_{c}(U, \mathcal{L}F(1))^{\vee}$(最后的 $\vee$ 表示 $\operatorname{Hom}(\cdot, \mathbb{Z}/n\mathbb{Z})$ 意义下的对偶).我们必须证明：由定理中的配对(把 $F$ 换成 $\mathcal{L}F(1)$)定义的映射
\[
\varphi^{r}(U, F) \colon H^{r}(U, F) \to T^{r}(U, F)
\]
对 $U$ 上所有有限局部常值层 $F$ 都是有限群的同构.

因为 $\operatorname{Hom}_{\mathbb{Z}/n\mathbb{Z}}(\cdot, \mathbb{Z}/n\mathbb{Z})$ 保持 $\mathbb{Z}/n\mathbb{Z}$-模序列的正合性，$F \mapsto \mathcal{L}F(1)$ 也如此(在茎上看).于是，有限局部常值层的正合序列
\[
0 \to F' \to F \to F'' \to 0
\]
产生正合序列
\[
\cdots \to T^{s}(U, F') \to T^{s}(U, F) \to T^{s}(U, F'') \to T^{s+1}(U, F') \to \cdots.
\]

第 1 步：设 $\varphi \colon U' \to U$ 是有限映射.定理对 $U'$ 上的 $F$ 成立，当且仅当它对 $U$ 上的 $\varphi_{*}F$ 成立.

\begin{proof}
回忆 $\varphi_{*}$ 正合且保持内射对象，因此对所有 $r$，$H^{r}(U, \varphi_{*}F) \cong H^{r}(U', F)$.由于对 $T^{s}$ 有类似的陈述，看到 $\varphi^{r}(U, \varphi_{*}F)$ 可以等同于 $\varphi^{r}(U', F)$.
\end{proof}

第 2 步：设 $V = U \setminus x$($x \in U$ 适当).则在标记为同构处有正合交换图
\[
\begin{tikzcd}
H^{r}_{x}(U, \mu_n) \arrow[r] \arrow[d] & H^{r}(U, \mu_n) \arrow[r] \arrow[d, "{\varphi^{r}(U, \mu_n)}"] & H^{r}(V, \mu_n) \arrow[r] \arrow[d, "{\varphi^{r}(V, \mu_n)}"] & H^{r+1}_{x}(U, \mu_n) \arrow[d] \\
H^{2-r}(x, \mathbb{Z}/n\mathbb{Z})^{\vee} \arrow[r] & H^{2-r}_{c}(U, \mathbb{Z}/n\mathbb{Z})^{\vee} \arrow[r] & H^{2-r}_{c}(V, \mathbb{Z}/n\mathbb{Z})^{\vee} \arrow[r] & H^{3-r}(x, \mathbb{Z}/n\mathbb{Z})^{\vee}
\end{tikzcd}
\]

\begin{proof}
上面的序列是空间对 $(U, V)$ 的正合序列.由 (\ref{prop:14.3})，当 $r = 2$ 时 $H^{r}_{x}(V, \mu_n) = H^{0}(x, \mathbb{Z}/n\mathbb{Z})$，否则为零.

下面的序列是
\[
0 \to j_{!}\mathbb{Z}/n\mathbb{Z} \to \mathbb{Z}/n\mathbb{Z} \to i_{*}i^{*}\mathbb{Z}/n\mathbb{Z} \to 0
\]
的紧支集上同调序列，其中 $j$ 与 $i$ 分别是 $V$ 与 $x$ 到 $U$ 的包含(见注记~\ref{rem:8.16}).

对 $r = 2$，未标记的竖映射是由配对
\[
H^{0}(x, \mathbb{Z}/n\mathbb{Z}) \times H^{0}(x, \mathbb{Z}/n\mathbb{Z}) \to H^{0}(x, \mathbb{Z}/n\mathbb{Z}) \cong \mathbb{Z}/n\mathbb{Z}
\]
(点 $x$ 的 Poincar\'e 对偶)给出的同构；对 $r \neq 2$，它是零群之间的映射.
\end{proof}

第 3 步：映射 $\varphi^{0}(U, \mathbb{Z}/n\mathbb{Z})$ 是有限群的同构.

\begin{proof}
在这种情况下，配对是
\[
H^{0}(U, \mathbb{Z}/n\mathbb{Z}) \times H^{2}_{c}(U, \mu_n) \to H^{2}_{c}(U, \mu_n) \cong \mathbb{Z}/n\mathbb{Z}.
\]
这里 $H^{0}(U, \mathbb{Z}/n\mathbb{Z}) = \mathbb{Z}/n\mathbb{Z}$，它在 $H^{2}_{c}(U, \mu_n)$ 上的作用由 $H^{2}_{c}(U, \mu_n)$ 的自然 $\mathbb{Z}/n\mathbb{Z}$-模结构定义.
\end{proof}

第 4 步：定理对 $r = 0$ 且 $F$ 局部常值成立.

\begin{proof}
设 $U' \to U$ 是使 $F|_{U'}$ 常值的有限平展覆盖.可以把 $F|_{U'}$ 嵌入 $U'$ 上的层 $F_0 = (\mathbb{Z}/n\mathbb{Z})^{s}$($s$ 适当).对 $F|_{U'} \hookrightarrow F_0$ 应用 $\varphi_{*}$，并把结果与自然包含 $F \hookrightarrow \varphi_{*}\varphi^{*}F$ 复合，得到序列
\[
0 \to F \to \varphi_{*}F_0 \to F'' \to 0
\]
中的第一个映射.余核 $F''$ 又局部常值(当 $U' \to U$ 选取为 Galois 时，$F''$ 在 $U'$ 上变成常值).考虑图
\[
\begin{tikzcd}
0 \arrow[r] & H^{0}(U, F) \arrow[r] \arrow[d, "{\varphi^{0}(U,F)}"] & H^{0}(U, \varphi_{*}F_0) \arrow[r] \arrow[d, "{\varphi^{0}(U,\varphi_{*}F_0)}"] & H^{0}(U, F'') \arrow[d, "{\varphi^{0}(U,F'')}"] \\
0 \arrow[r] & T^{0}(U, F) \arrow[r] & T^{0}(U, \varphi_{*}F_0) \arrow[r] & T^{0}(U, F'')
\end{tikzcd}
\]
``$0$''群实际上为零，但我们需要的一切只是箭头是同构.映射 $\varphi^{0}(U, \varphi_{*}F_0)$ 由第 1 步和第 3 步是同构.五引理表明 $\varphi^{0}(U, F)$ 是单射.由于这对所有局部常值层 $F$ 成立，$\varphi^{0}(U, F'')$ 是单射，五引理现在表明 $\varphi^{0}(U, F)$ 满射.最后，$H^{0}(U, F)$ 显然有限.
\end{proof}

第 5 步：映射 $\varphi^{1}(U, \mu_n)$ 是单射.

\begin{proof}
回忆 (\ref{ex:11.3}): $H^{1}(U, \mathbb{Z}/n\mathbb{Z}) = \operatorname{Hom}_{\mathrm{conts}}(\pi_1(U, \bar{u}), \mathbb{Z}/n\mathbb{Z})$.设 $s \in H^{1}(U, \mathbb{Z}/n\mathbb{Z})$，$\varphi \colon U' \to U$ 是对应于 $s$ 的核的 Galois 覆盖.则 $s$ 在 $H^{1}(U', \mathbb{Z}/n\mathbb{Z})$ 中映到零，后者同构于 $H^{1}(U, \varphi_{*}\mathbb{Z}/n\mathbb{Z})$.设 $F''$ 是 $\mathbb{Z}/n\mathbb{Z} \to \varphi_{*}(\mathbb{Z}/n\mathbb{Z})$ 的余核，考虑
\[
\begin{tikzcd}
H^{0}(U, \varphi_{*}\mathbb{Z}/n\mathbb{Z}) \arrow[r] \arrow[d] & H^{0}(U, F'') \arrow[r] \arrow[d] & H^{1}(U, \mathbb{Z}/n\mathbb{Z}) \arrow[r] \arrow[d, "{\varphi^{1}(U,\mathbb{Z}/n\mathbb{Z})}"] & H^{1}(U, \varphi_{*}\mathbb{Z}/n\mathbb{Z}) \arrow[d] \\
T^{0}(U, \varphi_{*}\mathbb{Z}/n\mathbb{Z}) \arrow[r] & T^{0}(U, F'') \arrow[r] & T^{1}(U, \mathbb{Z}/n\mathbb{Z}) \arrow[r] & T^{1}(U, \varphi_{*}\mathbb{Z}/n\mathbb{Z})
\end{tikzcd}
\]
由 $\varphi$ 的选取，$s$ 在 $H^{1}(U, \varphi_{*}\mathbb{Z}/n\mathbb{Z})$ 中映到零，而图追图表明：若 $s$ 也在 $T^{1}(U, \mathbb{Z}/n\mathbb{Z})$ 中映到零，则 $s$ 为零.
\end{proof}

第 6 步：映射 $\varphi^{r}(U, \mathbb{Z}/n\mathbb{Z})$ 对 $r = 1, 2$ 都是有限群的同构.

\begin{proof}
首先取 $U = X$.对 $r = 1$，$\varphi$ 是映射 $H^{1}(X, \mathbb{Z}/n\mathbb{Z}) \to H^{1}(X, \mu_n)^{\vee}$.因为 $\varphi^{1}(X, \mathbb{Z}/n\mathbb{Z})$ 单射(第 5 步)，且 $H^{1}(X, \mathbb{Z}/n\mathbb{Z})$ 有限，与 $H^{1}(X, \mu_n)$ 同阶(14.2, 14.4)，它是同构.对 $r = 2$，必须证明配对
\[
H^{2}(X, \mathbb{Z}/n\mathbb{Z}) \times \mu_n(k) \to H^{2}(X, \mu_n)
\]
完美，但这由 (\ref{rem:14.4}) 推出.

我们已经证明 $\varphi^{r}(X, \mathbb{Z}/n\mathbb{Z})$ 对 $r = 1, 2$ 是有限群的同构.要推出 $\varphi^{r}(U, \mathbb{Z}/n\mathbb{Z})$ 的同样论断，逐个去掉 $X \setminus U$ 的点，并把五引理应用于第 2 步的图.
\end{proof}

第 7 步：映射 $\varphi^{r}(U, F)$ 对 $r = 1, 2$ 且 $F$ 局部常值都是有限群的同构.

\begin{proof}
把第 4 步的论证应用两次.
\end{proof}

这完成了定理~\ref{thm:14.7} 证明的梗概.

\subsection*{Hochschild-Serre 谱序列}

本小节中，$X$ 是任意簇(或概形).

\begin{theorem}\label{thm:14.9}
设 $\varphi \colon Y \to X$ 是带 Galois 群 $G$ 的 Galois 覆盖.对 $X_{\mathrm{et}}$ 上的任意层 $F$，存在谱序列
\begin{equation}
H^{r}(G, H^{s}(Y_{\mathrm{et}}, F|_{Y})) \Rightarrow H^{r+s}(X_{\mathrm{et}}, F). \tag{1}
\end{equation}
\end{theorem}

\begin{proof}
对 $X$ 上的任意层 $F$，$F(Y)$ 是左 $G$-模，且 $F(X) = F(Y)^{G}$(见命题~\ref{prop:6.4}).因此函子
\[
\mathrm{Sh}(X_{\mathrm{et}}) \xrightarrow{F \mapsto F(Y)} G\text{-}\mathrm{Mod} \xrightarrow{M \mapsto M^{G}} \mathbf{Ab}
\]
的复合是 $\Gamma(X, \cdot)$.一旦证明 $\mathcal{I}$ 内射蕴含对 $r > 0$，$H^{r}(Y, \mathcal{I}(Y)) = 0$，定理就由 (\ref{thm:10.7}) 推出.但若 $\mathcal{I}$ 作为层内射，则它作为预层也内射，从而对 $r > 0$，$\check{H}^{r}(Y/X, \mathcal{I}) = 0$(见 \S 10).如我们在 (\ref{ex:10.1}) 中所述，$\check{H}^{r}(Y/X, \mathcal{I}) = \check{H}^{r}(G, \mathcal{I}(Y))$.
\end{proof}

\begin{quote}
\textbf{注记} 当 $Y$ 与 $X$ 是域的谱时，(1) 变成 Hochschild-Serre 谱序列，因此它也通常被称为 Hochschild-Serre 谱序列.然而，正如 Serre 所解释的，\footnote{``En septembre 1950, un peu avant de clarifier ce qui allait \^etre ma th\`ese, j'ai vu que les suites spectrales de Leray s'appliquaient aussi \`a la cohomologie des extensions de groupes. C'\'etait d'ailleurs presque \'evident vu ce qu'avaient d\'ej\`a fait Cartan-Leray: ils avaient montr\'e que, si $X \to Y$ est un rev\^etement galoisien de groupe $\Gamma$, il y a une suite spectrale qui part de $H^{*}(\Gamma, H^{*}(X))$ et aboutit \`a $H^{*}(Y)$. [Cette suite spectrale est souvent appel\'ee maintenant ``de Hochschild-Serre''---je n'y peux rien].'' Serre, Notices AMS, Sept. 2011, p.1087.}``Cartan-Leray 谱序列''会更合适.
\end{quote}

\begin{remark}\label{rem:14.10}
设
\[
\cdots \to Y_i \to Y_{i-1} \to \cdots \to Y_1 \to Y_0 = X
\]
是其中每个映射 $Y_i \to Y_j$ 都 Galois 的塔.设 $G_i$ 是 $Y_i$ 在 $X$ 上的 Galois 群，$G = \varprojlim G_i$.对每个 $i$，存在谱序列
\[
H^{r}(G_i, H^{s}(Y_i, F_i)) \Rightarrow H^{r+s}(X, F), \qquad F_i = F|_{Y_i}.
\]
取逆向极限，得到谱序列
\[
H^{r}(G, H^{s}(Y_\infty, F_\infty)) \Rightarrow H^{r+s}(X, F).
\]
这里 $H^{r}(G, \cdot)$ 是用连续上链计算的 pro-有限群 $G$ 的上同调群，$Y_\infty = \varprojlim Y_i$，$F_\infty$ 是 $F$ 在 $Y_\infty$ 上的逆像(比较 10.9 与 CFT II 4.4).
\end{remark}

\begin{example}\label{ex:14.11}
设 $A$ 是分式域为 $K$ 的 Dedekind 整环，$K^{\mathrm{sep}}$ 是 $K$ 的可分闭包.设 $K^{\mathrm{un}}$ 是所有这样的子域 $L \subset K^{\mathrm{sep}}$ 的复合：$[L \colon K] < \infty$，且 $A$ 在 $L$ 中的整闭包在所有素理想上都在 $A$ 上非分歧.设 $G = \operatorname{Gal}(K^{\mathrm{un}}/K)$---它是 $\pi_1(U, \bar{\xi})$，其中 $U = \operatorname{Spec} A$，$\bar{\xi}$ 是几何点 $\operatorname{Spec} K^{\mathrm{sep}} \to U$.设 $F_M$ 是对应于离散 $\pi_1$-模 $M$ 的 $U$ 上的局部常值层.则存在谱序列
\[
H^{r}(\pi_1(U, \bar{\xi}), H^{s}(\widetilde{U}, F_M)) \Rightarrow H^{r+s}(U, F_M).
\]
若对 $s > 0$，$H^{s}(\widetilde{U}, F_M) = 0$，则这给出对所有 $r$ 的同构
\[
H^{r}(\pi_1(U, \bar{\xi}), M) \cong H^{r}(U, F_M).
\]
\end{example}

\subsection*{仿射曲线上局部常值层的上同调}

本节中，我们假定：若 $\operatorname{char}(k) = p \neq 0$，则层没有 $p$-挠.

设 $X$ 是域 $k$ 上的完备连通非奇异曲线.对 $X$ 的任何非空有限闭点集 $S$，$U = X \setminus S$ 仿射.(Riemann-Roch 定理为我们提供 $U$ 上的非常数函数 $f$，即非常数函数 $f \colon U \to \mathbb{A}^{1}$.)

\begin{proposition}\label{prop:14.12}
设 $U$ 是代数闭域上的非奇异仿射曲线.对 $U$ 上的任何局部常值挠层 $F$，对 $r \geq 2$，$H^{r}(U_{\mathrm{et}}, F) = 0$.
\end{proposition}

\begin{proof}
首先证明 $H^{2}(U, \mu_n) = 0$.由 Weil 因子序列，我们发现 $H^{1}(U, \mathbb{G}_{m}) = \operatorname{Pic}(U)$，且对 $r > 1$，$H^{r}(U, \mathbb{G}_{m}) = 0$，因此 Kummer 序列表明 $H^{2}(U, \mu_n) = \operatorname{Pic}(U)/n\operatorname{Pic}(U)$，且对 $r > 2$，$H^{r}(U, \mu_n) = 0$.

设 $X$ 是包含 $U$ 的完备非奇异曲线.由于 $U$ 至少省略 $X$ 的一个点，投影映射
\[
\mathcal{D}iv(X) \to \mathcal{D}iv(U)
\]
\[
\bigoplus_{x \in X} \mathbb{Z} \to \bigoplus_{x \in U} \mathbb{Z}
\]
把 $\mathcal{D}iv^{0}(X) \overset{\mathrm{def}}{=} \{\sum n_x [x] \mid \sum n_x = 0\}$ 映到 $\mathcal{D}iv(U)$ 上.由于投影映射把主因子送到主因子，看到 $\operatorname{Pic}(U)$ 是 $\operatorname{Pic}^{0}(X)$ 的商，而后者是可除的.因此 $\operatorname{Pic}(U)$ 也可除，从而 $H^{2}(U, \mu_n) = 0$.

其次，设 $F$ 是 $U$ 上带有限纤维的局部常值层.存在有限 Galois 覆盖 $\varphi \colon U' \to U$，使得 $F|_{U'}$ 常值，并有满射``迹映射'' $\varphi_{*}\varphi^{*}F \to F$.因为对 $r \geq 2$，$H^{r}(U, \varphi_{*}\varphi^{*}F) = 0$，由此推出对 $r \geq 2$，$H^{r}(U, F) = 0$.(设 $G$ 是 $U'$ 在 $U$ 上的 Galois 群.对任何平展映射 $V \to U$，$\Gamma(V, \varphi_{*}\varphi^{*}F) = \Gamma(U' \times_{U} V, F)$，它是带固定模 $\Gamma(V, F)$ 的 $G$-模.迹映射是 $s \mapsto \sum_{g \in G} gs$.要看出它满射，在茎上看.)

最后，局部常值挠层是带有限纤维的局部常值层的正向极限.
\end{proof}

\begin{proposition}\label{prop:14.13}
设 $U$ 是带 $\pi_1(U, \bar{u}) = 1$(对某个(从而每个)几何点 $\bar{u} \to U$)的连通簇(或概形).则对 $U$ 上的任何局部常值层，$H^{1}(U, F) = 0$.
\end{proposition}

\begin{proof}
由 (\ref{prop:6.16})，任何局部常值层都常值，而由 (\ref{ex:11.3}), $H^{1}(U, F_M) = \operatorname{Hom}(\pi_1(U, \bar{u}), M) = 0$.
\end{proof}

\begin{theorem}\label{thm:14.14}
设 $U$ 是代数闭域 $k$ 上的光滑仿射曲线，$F$ 是对应于离散 $\pi_1(U, \bar{u})$-模 $M$ 的 $U$ 上的局部常值挠层.则对所有 $r$，$H^{r}(U, F) \cong H^{r}(\pi_1(U, \bar{u}), M)$.
\end{theorem}

\begin{proof}
设 $\widetilde{U}$ 是 $U$ 的通用覆盖簇.则由 (\ref{prop:14.12})，对 $r \geq 2$，$H^{r}(\widetilde{U}, F) = 0$，且因为 $\widetilde{U}$ 单连通，$H^{1}(\widetilde{U}, F) = 0$.现在可以应用 (\ref{ex:14.11}).
\end{proof}

\subsection*{有限域上的曲线}

设 $k$ 是有限域，$\Gamma = \operatorname{Gal}(k^{\mathrm{al}}/k)$.则 $\Gamma \cong \widehat{\mathbb{Z}}$，其中 $1 \in \widehat{\mathbb{Z}}$ 对应于 $\Gamma$ 中的 Frobenius 元素 $F$.本节中，$n$ 是与 $k$ 的特征互素的整数.

有典范同构
\[
H^{1}(\Gamma, \mathbb{Z}/n\mathbb{Z}) \to \operatorname{Hom}_{\mathrm{conts}}(\Gamma, \mathbb{Z}/n\mathbb{Z}) \xrightarrow{\varphi \mapsto \varphi(F)} \mathbb{Z}/n\mathbb{Z}.
\]
此外，对 $r > 1$ 与任何挠离散 $\Gamma$-模(见~\ref{item:13.5} 与命题~\ref{prop:13.6}), $H^{r}(\Gamma, M) = 0$.

\begin{proposition}\label{prop:14.15}
设 $M$ 是有限离散 $\Gamma$-模，$\check{M} = \operatorname{Hom}(M, \mathbb{Z}/n\mathbb{Z})$(对偶 Abel 群).则杯积配对
\[
H^{r}(\Gamma, M) \times H^{1-r}(\Gamma, \check{M}) \to H^{1}(\Gamma, \mathbb{Z}/n\mathbb{Z}) \cong \mathbb{Z}/n\mathbb{Z}
\]
是有限群的完美配对.
\end{proposition}

\begin{proof}
按定义，$H^{0}(\Gamma, M) = M^{\Gamma}$，即 $\Gamma$ 平凡作用的最大的 $M$ 的子群.把交叉同态 $f \colon \Gamma \to M$ 送到它在 Frobenius 元素上的值的映射诱导 $H^{1}(\Gamma, M)$ 到 $M_{\Gamma}$(即 $\Gamma$ 平凡作用的最大的 $M$ 的商)的同构.在对偶 $M \leftrightarrow \check{M}$ 中，$\Gamma$ 平凡作用的最大的子群对应于 $\Gamma$ 平凡作用的最大的商群，即完美配对
\[
M \times \check{M} \to \mathbb{Z}/n\mathbb{Z}
\]
诱导完美配对
\[
M^{\Gamma} \times (\check{M})_{\Gamma} \to \mathbb{Z}/n\mathbb{Z}.
\]
\end{proof}

该命题可以更简洁地陈述为：对任何被 $n$ 消灭的有限 $\Gamma$-模 $M$，
\[
(M^{\Gamma})^{\vee} = (M^{\vee})_{\Gamma}, \qquad (M_{\Gamma})^{\vee} = (M^{\vee})^{\Gamma}.
\]

设 $U$ 是有限域 $k$ 上的连通非奇异曲线，$\bar{U}$ 是由基变换得到的 $k^{\mathrm{al}}$ 上的曲线.对任何挠 $\Gamma$-模 $M$，
\[
H^{r}(\Gamma, M) = M^{\Gamma}, M_{\Gamma}, 0 \quad \text{分别地, 对 } r = 0, 1, \geq 2.
\]
因此，对 $U$ 上的任何挠层 $F$，Hochschild-Serre 谱序列给出短正合序列
\begin{equation}
0 \to H^{r-1}(\bar{U}, F)_{\Gamma} \to H^{r}(U, F) \to H^{r}(\bar{U}, F)^{\Gamma} \to 0, \tag{2}
\end{equation}
对紧支集上同调群类似地，
\[
0 \to H^{r-1}_{c}(\bar{U}, F)_{\Gamma} \to H^{r}_{c}(U, F) \to H^{r}_{c}(\bar{U}, F)^{\Gamma} \to 0.
\]
若 $F$ 是被 $n$ 消灭的 $U$ 上的有限局部常值层，则该序列中的群有限，而按命题~\ref{prop:14.15}，它的对偶是
\[
0 \to (H^{r}_{c}(\bar{U}, F)^{\vee})^{\Gamma} \to H^{r}_{c}(U, F)^{\vee} \to (H^{r-1}_{c}(\bar{U}, F)^{\vee})_{\Gamma} \to 0.
\]
用定理~\ref{thm:14.7} 替换该序列中的两个群，得到序列
\begin{equation}
0 \to H^{2-r}(\bar{U}, \mathcal{L}F(1))^{\Gamma} \to H^{r}_{c}(U, F)^{\vee} \to H^{3-r}(\bar{U}, \mathcal{L}F(1))_{\Gamma} \to 0. \tag{3}
\end{equation}
比较 (2) 与 (3) 提示：

\begin{theorem}\label{thm:14.16}[Poincar\'e 对偶]
对 $U$ 上的任何有限局部常值层 $F$ 与整数 $r \geq 0$，有有限群的典范完美配对
\[
H^{r}_{c}(U, F) \times H^{3-r}(U, \mathcal{L}F(1)) \to H^{3}_{c}(U, \mu_n) \cong \mathbb{Z}/n\mathbb{Z}.
\]
\end{theorem}

\begin{proof}
唯一的困难是定义配对，并验证它与 $\bar{U}$ 上的配对相容，即证明
\[
\begin{tikzcd}
0 \arrow[r] & H^{r-1}_{c}(\bar{U}, F)_{\Gamma} \arrow[r] \arrow[d] & H^{r}_{c}(U, F) \arrow[r] \arrow[d] & H^{r}_{c}(\bar{U}, F)^{\Gamma} \arrow[r] \arrow[d] & 0 \\
0 \arrow[r] & (H^{3-r}(\bar{U}, \mathcal{L}F(1))^{\Gamma})^{\vee} \arrow[r] & H^{3-r}(U, \mathcal{L}F(1))^{\vee} \arrow[r] & (H^{2-r}(\bar{U}, \mathcal{L}F(1))_{\Gamma})^{\vee} \arrow[r] & 0
\end{tikzcd}
\]
交换.
\end{proof}

\subsection*{有限域上函数域的 Galois 上同调中的正合序列}

设 $X$ 是有限域 $k$ 上的完备非奇异连通曲线，$U$ 是 $X$ 的开子集.设 $F$ 是对应于阶与 $k$ 的特征互素的有限 $\pi_1(U, \bar{u})$-模 $M$ 的 $U$ 上的层.我们希望纯粹用 Galois 上同调群解释正合序列
\[
\cdots \to H^{r}_{Z}(X, j_{!}F) \to H^{r}(X, j_{!}F) \to H^{r}(U, F) \to \cdots.
\]
设 $G = \pi_1(U, \bar{u})$，对一切 $v \in X \setminus U$，设 $G_v$ 是 $v$ 处的分解群.则
\[
H^{r}(U, F) \cong H^{r}(G, M)
\]
\[
H^{r}(X, j_{!}F) \overset{\mathrm{def}}{=} H^{r}_{c}(U, F) \cong H^{3-r}(U, \mathcal{L}F(1))^{\vee} \cong H^{3-r}(G, \check{M}(1))^{\vee}.
\]
还剩下计算 $H^{r}_{Z}(X, j_{!}F)$.由切除，$H^{r}_{Z}(X, j_{!}F) = \bigoplus H^{r}_{z}(X_z, j_{!}F)$，其中 $X_z = \operatorname{Spec} \mathcal{O}^{\mathrm{h}}_{X, z}$.

\begin{lemma}\label{lem:14.17}
设 $V = \operatorname{Spec} A$，其中 $A$ 是离散赋值环，$u$ 与 $z$ 分别是 $V$ 的开点与闭点.对 $u$ 上的任何层 $F$ 与所有 $r$，$H^{r}(V, j_{!}F) = 0$，因此空间对 $(V, u)$ 的正合序列中的边缘映射
\[
H^{r-1}(u, F) \to H^{r}_{z}(V, F)
\]
对所有 $r$ 都是同构.
\end{lemma}

\begin{proof}
见我的书 \emph{Arithmetic Duality Theorems}, II 1.1.
\end{proof}

于是
\[
H^{r}_{Z}(U, j_{!}F) = \bigoplus_{v \in Z} H^{r-1}(G_v, M).
\]

\begin{theorem}\label{thm:14.18}
对阶与 $k$ 的特征互素的任何有限离散 $G$-模 $M$，有典范正合序列
\[
\begin{array}{c}
0 \to H^{0}(G, M) \to \displaystyle\bigoplus_{v \in Z} H^{0}(G_v, M) \to H^{2}(G, \check{M}(1))^{\vee} \\
\downarrow \\
H^{1}(G, \check{M}(1))^{\vee} \to \displaystyle\bigoplus_{v \in Z} H^{1}(G_v, M) \to H^{1}(G, M) \\
\downarrow \\
H^{2}(G, M) \to \displaystyle\bigoplus_{v \in Z} H^{2}(G_v, M) \to H^{0}(G, \check{M}(1))^{\vee} \to 0
\end{array}
\]
\end{theorem}

\begin{proof}
在空间对 $(X, U)$ 的 $j_{!}F$ 的正合序列中，对 $F$($U$ 上对应于 $M$ 的局部常值层)作上面描述的代换.
\end{proof}

\subsection*{数域的 Galois 上同调中的正合序列}

设 $K$ 是数域，$S$ 是 $K$ 的有限素点集，包含无限素点.设 $R_S$ 是 $K$ 中的 $S$-整数环，即 $K$ 中满足 $\operatorname{ord}_v(a) \geq 0$($v \notin S$)的元素 $a$ 组成的环.

\begin{theorem}\label{thm:14.19}
设 $G_S$ 是 $K$ 在 $S$ 外非分歧的最大扩张的 Galois 群，$M$ 是其阶在 $R_S$ 中为单位的有限离散 $G_S$-模.有典范正合序列
\[
\begin{array}{c}
0 \to H^{0}(G_S, M) \to \displaystyle\bigoplus_{v \in Z} H^{0}(G_v, M) \to H^{2}(G_S, \check{M}(1))^{\vee} \\
\downarrow \\
H^{1}(G_S, \check{M}(1))^{\vee} \to \displaystyle\bigoplus_{v \in Z} H^{1}(G_v, M) \to H^{1}(G_S, M) \\
\downarrow \\
H^{2}(G_S, M) \to \displaystyle\bigoplus_{v \in Z} H^{2}(G_v, M) \to H^{0}(G_S, \check{M}(1))^{\vee} \to 0
\end{array}
\]
\end{theorem}

定理~\ref{thm:14.19} 是 Tate 的非常重要结果---见他在 ICM 1962 的报告(Poitou 大约同时得到类似结果).它的直接证明(基于 Tate 的未发表证明)见我的书 \emph{Arithmetic Duality Theorems} 的 I.4.或者，它可以由 $R_S$ 的平展上同调的一个定理恢复(见下面 (\ref{thm:14.22})).

\subsection*{配对与 Ext 群}

本小节解释如何定义上述对偶定理中的配对，以及如何把定理推广到更一般的层.

首先陈述推广 (\ref{thm:14.7}) 的定理，然后解释它.

\begin{theorem}\label{thm:14.20}
设 $U$ 是代数闭域 $k$ 上的非奇异曲线.对 $U$ 上 $\mathbb{Z}/n\mathbb{Z}$-模的一切可构造层 $F$ 与所有 $r \geq 0$，有有限群的典范完美配对
\[
H^{r}_{c}(U, F) \times \operatorname{Ext}^{2-r}_{U, n}(F, \mu_n) \to H^{r}_{c}(U, \mu_n) \cong \mathbb{Z}/n\mathbb{Z}.
\]
\end{theorem}

\textbf{可构造层} 域 $k$ 上的曲线 $U$ 上的层 $F$ 称为可构造的，如果
\begin{enumerate}
  \item 对 $U$ 的某个非空开子集 $V$，$F|_{V}$ 局部常值；
  \item $F$ 的茎有限.
\end{enumerate}
应用 (\ref{prop:8.17}) 与 (\ref{prop:6.16})，看到给出 $U$ 上的可构造层必须给出
\begin{enumerate}
  \item $U$ 的有限闭点集 $S$；设 $V = U \setminus S$；
  \item 有限离散 $\pi_1(V, \bar{\xi})$-模 $M_V$(这里 $\xi$ 是 $U$ 的一般点)；
  \item 对每个 $v \in S$，有限离散 $\operatorname{Gal}(k(v)^{\mathrm{sep}}/k(v))$-模 $M_v$；
  \item 对群的共变同态 $M_v \to M_V$.
\end{enumerate}

Abel 群层 $F$ 可以看作 $\mathbb{Z}$-模层(这里 $\mathbb{Z}$ 是常值层).说它是 $\mathbb{Z}/n\mathbb{Z}$-模层无非是说 $nF = 0$.

\textbf{Ext 群} 对任何簇(或概形) $X$，设 $\mathrm{Sh}(X_{\mathrm{et}, n})$ 是 $X_{\mathrm{et}}$ 上 $\mathbb{Z}/n\mathbb{Z}$-模层的范畴.与全范畴一样，$\mathrm{Sh}(X_{\mathrm{et}, n})$ 有足够的内射对象.对 $\mathrm{Sh}(X_{\mathrm{et}, n})$ 中的层 $F$，无论用 $\mathrm{Sh}(X_{\mathrm{et}})$ 还是 $\mathrm{Sh}(X_{\mathrm{et}, n})$ 中的内射分解计算，$H^{r}(X_{\mathrm{et}}, F)$ 都相同(见 EC III 2.25).然而，若 $F$ 与 $G$ 是两个被 $n$ 消灭的层，则 $\operatorname{Ext}^{r}(F, G)$ 依赖于在哪个范畴中计算.当 $X$ 是代数闭域的谱时已经可以看到这一点.这里 $\mathrm{Sh}(X_{\mathrm{et}})$ 是 $\mathbb{Z}$-模的范畴，$\mathrm{Sh}(X_{\mathrm{et}, n})$ 是 $\mathbb{Z}/n\mathbb{Z}$-模的范畴.在第一个范畴中，$\operatorname{Ext}^{1}(\mathbb{Z}/n\mathbb{Z}, \mathbb{Z}/n\mathbb{Z}) \neq 0$，因为
\[
0 \to \mathbb{Z}/n\mathbb{Z} \to \mathbb{Z}/n^{2}\mathbb{Z} \xrightarrow{n} \mathbb{Z}/n\mathbb{Z} \to 0
\]
不分裂，而在第二个范畴中 $\operatorname{Ext}^{1}(\mathbb{Z}/n\mathbb{Z}, \mathbb{Z}/n\mathbb{Z}) = 0$，因为 $\mathbb{Z}/n\mathbb{Z}$ 内射(作为 $\mathbb{Z}/n\mathbb{Z}$-模).

$\operatorname{Ext}^{r}_{U, n}(F, G)$ 指在范畴 $\mathrm{Sh}(U_{\mathrm{et}, n})$ 中计算的 Ext 群.

\begin{lemma}\label{lem:14.21}
若 $F$ 局部常值，则 $\operatorname{Ext}^{r}_{U, n}(F, \mu_n) \cong H^{r}(U_{\mathrm{et}}, \mathcal{L}F(1))$，其中 $\mathcal{L}F(1)$ 是层
\[
V \mapsto \operatorname{Hom}_{V}(F|_{V}, \mu_n).
\]
\end{lemma}

\begin{proof}
(证明的梗概.) 对任意簇(或概形) $Y$ 上的层 $F$ 与 $G$，定义 $\mathcal{H}om(F, G)$ 为层
\[
V \mapsto \operatorname{Hom}_{V_{\mathrm{et}}}(F|_{V}, G|_{V})
\]
---容易检验 $\mathcal{H}om(F, G)$ 确实是层.设 $F_0$ 是带有限茎的 $\mathbb{Z}/n\mathbb{Z}$-模局部常值层.函子
\[
G \mapsto \mathcal{H}om(F_0, G) \colon \mathrm{Sh}(Y_{\mathrm{et}, n}) \to \mathrm{Sh}(Y_{\mathrm{et}, n})
\]
左正合，因此可以构造它的右导出函子 $\mathcal{E}xt^{r}(F_0, \cdot)$.它们称为局部 Ext 群.它们与整体 Ext 群由谱序列
\[
H^{r}(Y_{\mathrm{et}}, \mathcal{E}xt^{s}(F_0, G)) \Rightarrow \operatorname{Ext}^{s}(F_0, G)
\]
联系.在 $F_0$ 的假设下，茎
\[
\mathcal{E}xt^{r}(F_0, G)_{\bar{y}} = \operatorname{Ext}^{r}((F_0)_{\bar{y}}, G_{\bar{y}})
\]
对所有几何点 $\bar{y}$ 成立.右边的 Ext 在 $\mathbb{Z}/n\mathbb{Z}$-模范畴中形成.因此，在引理的情形中，
\[
\mathcal{E}xt^{r}(F, \mu_n)_{\bar{y}} = \operatorname{Ext}^{r}(F_{\bar{y}}, \mu_n) \cong \operatorname{Ext}^{r}(F_{\bar{y}}, \mathbb{Z}/n\mathbb{Z}) = 0 \quad (r > 0),
\]
而谱序列给出同构
\[
H^{r}(X_{\mathrm{et}}, \mathcal{H}om(F, \mu_n)) \cong \operatorname{Ext}^{r}_{X, n}(F, \mu_n),
\]
正如所要求的.更多细节见 EC II 3.20, III 1.22, III 1.31.
\end{proof}

于是，对局部常值层 $F$，定理~\ref{thm:14.20} 中的配对具有形式
\[
H^{r}_{c}(U, F) \times H^{2-r}(U, \mathcal{L}F(1)) \to H^{2}_{c}(U, \mu_n) \cong \mathbb{Z}/n\mathbb{Z}.
\]

\textbf{配对} 定义配对有各种方法，但(幸运的是)它们都一致，至少到符号.也许最简单的是把 Ext 群解释为扩张的等价类，然后把配对定义为迭代边缘映射.

在任何 Abel 范畴中，$\operatorname{Ext}^{r}(A, B)$($r > 0$)的元素可以解释为 $r$ 重扩张
\[
0 \to B \to E_1 \to \cdots \to E_r \to A \to 0
\]
的等价类.有自然配对
\[
\operatorname{Ext}^{r}(A, B) \times \operatorname{Ext}^{s}(B, C) \to \operatorname{Ext}^{r+s}(A, C),
\]
它既可以通过拼接扩张定义，也可以作为迭代边缘映射定义，即把扩张
\[
0 \to C \to E_1 \to \cdots \to E_s \to B \to 0
\]
拆成短正合序列
\[
\begin{array}{c}
0 \to K_s \to E_s \to B \to 0\\
0 \to K_{s-1} \to E_{s-1} \to K_s \to 0\\
\vdots
\end{array}
\]
并形成边缘映射
\[
\operatorname{Ext}^{r}(A, B) \to \operatorname{Ext}^{r+1}(A, K_s) \to \operatorname{Ext}^{r+2}(A, K_{s-1}) \to \cdots \to \operatorname{Ext}^{r+s}(A, C).
\]
见 Mitchell, B., \emph{Theory of Categories}, Academic Press, 1965, VII.3 或 Bourbaki, N., \emph{Alg\`ebre}, Chap. X.

在我们的情形中，$\operatorname{Ext}^{s}(F, \mu_n)$ 的元素由正合序列
\[
0 \to \mu_n \to E_1 \to \cdots \to E_s \to F \to 0
\]
表示.应用 $j_{!}$ 得到扩张
\[
0 \to j_{!}\mu_n \to j_{!}E_1 \to \cdots \to j_{!}E_s \to j_{!}F \to 0,
\]
它定义迭代边缘映射
\[
H^{r}(X, j_{!}F) \to H^{r+s}(X, j_{!}\mu_n).
\]

\textbf{定理~\ref{thm:14.20} 的证明} 对可构造层 $F$，对 $r \geq 2$，$\operatorname{Ext}^{r}_{X, n}(F, \cdot) = 0$，且 $\operatorname{Ext}^{1}_{X, n}(F, \cdot)$ 的支集在有限集($F|_{V}$ 局部常值的开集 $V$ 的补集)上，因此 $\operatorname{Ext}^{r}_{X, n}(F, \mu_n)$ 的消失由第一个论断与 (\ref{lem:14.21}) 证明中的谱序列推出.定理~\ref{thm:14.7} 在 $F$ 局部常值的情形蕴含定理~\ref{thm:14.20}，五引理论证然后允许得到完整定理(见 EC p.177，第 2 步).

\subsection*{数域中的整数环}

如前，设 $K$ 是全虚数域，$S$ 是 $K$ 的包含无限素点的有限素点集.设 $R_S$ 是 $K$ 中的 $S$-整数环，$n$ 是在 $R_S$ 中为单位的整数.

\begin{theorem}\label{thm:14.22}
设 $U = \operatorname{Spec} R_S$.

(a) 有典范同构 $H^{3}_{c}(U, \mu_n) \cong \mathbb{Z}/n\mathbb{Z}$.

(b) 对 $U$ 上 $\mathbb{Z}/n\mathbb{Z}$-模的任何可构造层 $F$，有有限群的典范完美配对
\[
H^{r}_{c}(U, F) \times \operatorname{Ext}^{3-r}_{U, n}(F, \mu_n) \to H^{3}_{c}(U, \mu_n) \cong \mathbb{Q}/\mathbb{Z}.
\]
\end{theorem}

该定理的证明见我的 \emph{Arithmetic Duality Theorems}, II 3.或者，在 $F$ 局部常值的情形，定理~\ref{thm:14.19} 蕴含定理~\ref{thm:14.22}，而由此容易推出一般情形(同上，II 3.5).
